% ==============================================================================
% CHAPTER 1: INTRODUCTION
% ==============================================================================

\chapter{Introduction}
\label{chap:introduction}

Modern IT infrastructures have undergone a deep transformation over the last decade. The rapid shift toward cloud-native, containerized microservices has decomposed monolithic applications into dozens of cooperating processes distributed across heterogeneous hosts, and the digital attack surface has expanded accordingly. In this landscape, cyber threats have grown in both sophistication and amount: attackers chain exploits across services, automate attacks and weaponize compromised components faster than human operators can respond. The sheer volume of sensitive data housed in these environments, and the regulatory cost of exposing it, have turned the ability to \emph{defend autonomously} from a convenience into a necessity. Within this context, the autonomic computing paradigm has become a cornerstone of modern security engineering, promising to reduce human operational burden and to compress the time between detection and response. However, while automation resolves many latency and coordination challenges, it introduces a new and subtler systemic vulnerability: if the automated security logic itself becomes the target, a single compromise can be turned into a vector for devastating privilege escalation and system-wide failure. This thesis investigates how to build the decision substrate of such an autonomous defender.

% : : : : : : : : : : : : : : : : : : : : : : : : : : 
\section{The Paradigm of Self-Protecting Systems}
\label{sec:intro_paradigm}
% : : : : : : : : : : : : : : : : : : : : : : : : : : 

Self-Protecting Systems \cite{10.1145/2555611} (SPS) are autonomic software systems designed to detect, analyze, and mitigate security threats dynamically with limited or no human intervention. They are one of the four ``self-*'' facets: alongside self-configuration, self-optimization, and self-healing: through which IBM's autonomic computing manifesto \cite{kephart2003vision} originally characterized self-managing software, and they operationalize the self-protection objective through a continuous feedback loop.

Following the classic autonomic feedback loop, an SPS typically comprises two primary macro-components: a \textit{managed system} and an \textit{autonomic manager}. The managed system represents the production environment, the services, containers, or software applications that are directly exposed to potential attacks. Conversely, the autonomic manager serves as the central controller responsible for executing the security loop, and its internal logic is structured around the MAPE-K cycle \cite{kephart2003vision, 7194653}: it \emph{monitors} the managed system through sensors, \emph{analyzes} the collected observations, \emph{plans} a response consistent with high-level policies, and \emph{executes} it through effectors, all grounded on a shared \emph{knowledge} base.

To achieve continuous protection, the autonomic manager organizes its operations into a pipeline governed by two main functional sub-components:
\begin{itemize}
    \item \textbf{Intrusion Detection Systems (IDS):} Passive or active security sensors distributed across the managed system to monitor traffic, logs, or system calls and detect anomalous behavior or known attack signatures.
    \item \textbf{Intrusion Response Systems (IRS):} Decision-making engines responsible for selecting and executing the appropriate response actions (countermeasures) to neutralize or mitigate the detected threats.
\end{itemize}

In a self-protecting loop, the quality of protection is therefore bounded by two coupled properties: the correctness and diversity of the observations that the IDS plane produces, and the timeliness and integrity with which the IRS plane turns those observations into enforced countermeasures. The work of this thesis concentrates on the latter: not on how attacks are detected, but on whether the decision substrate that sits between detection and actuation can be trusted to remain correct, responsive, and available even when the infrastructure around it is under active attack.

% : : : : : : : : : : : : : : : : : : : : : : : : : : 
\section{The Vulnerability of the Autonomic Manager}
\label{sec:intro_vulnerability}
% : : : : : : : : : : : : : : : : : : : : : : : : : : 

Although highly effective in mitigating threats directed at the managed system, this canonical design suffers from a critical structural vulnerability: the centralization of the autonomic manager. An adversary who penetrates the infrastructure can target the autonomic manager itself. By compromising this central controller, the attacker can silence intrusion alerts, disable response mechanisms, or even manipulate the manager into executing malicious countermeasures (e.g., disconnecting legitimate network segments or isolating critical database services).

This is the ``Who watches the watchers?'' dilemma of autonomic computing \cite{chess2003security, pekaric2023systematic}: the very component that is supposed to enforce security becomes a high-value target, and protecting it through conventional defense-in-depth only hardens a single point of failure rather than removing it. If the security logic is centralized and compromised, the integrity of the entire infrastructure is undermined, rendering the self-protecting features not only ineffective but actively weaponized against the system. The state of the art has explored several responses to this dilemma, each with its own limitations. Traditional hardening of a centralized manager \cite{pekaric2023systematic} reduces the attack surface but does not eliminate the single point of failure. Trusted Execution Environments (TEEs), and in particular Intel SGX, have established themselves as the current benchmark by executing the manager's critical logic inside a hardware enclave; yet they remain exposed to microarchitectural side-channel attacks such as Spectre, Meltdown, and ZombieLoad, they introduce vendor lock-in, and, crucially, they do not solve the availability problem, since an attacker can simply crash or disconnect the single host running the enclave. A structural solution is therefore needed; one that removes the centralization assumption altogether rather than reinforcing it.

% : : : : : : : : : : : : : : : : : : : : : : : : : : 
\section{Blockchain as a Mechanism of Trust}
\label{sec:intro_blockchain}
% : : : : : : : : : : : : : : : : : : : : : : : : : : 

To eliminate the single point of failure represented by a centralized manager, our research \cite{11126595} has proposed distributing the autonomic manager across a permissioned blockchain network. In this decentralized architecture, blockchain technology acts as a tamper-proof ledger and a distributed state machine (DSM) that logs security events and executes response logic securely. Replicating the manager's state across multiple independent nodes, and encoding the response logic as a deterministic transition function over a globally agreed state, converts the manager from an isolated, hardening-dependent process into a resilient, physically separated network whose correctness rests on a formal fault model rather than on the impregnability of a single host.

This distributed approach establishes fundamental security guarantees that collectively safeguard the autonomic loop. To prevent a single compromised IDS sensor from injecting false alerts and triggering unauthorized countermeasures, the state machine enforces a consensual voting mechanism. Under this scheme, updates to the shared security state are only committed once they receive corroborating reports from multiple independent, heterogeneous IDSs, ensuring that isolated sensor compromises do not lead to malicious actuations.

Also, by deploying the autonomic manager across a network of validator nodes running a Byzantine Fault-Tolerant (BFT) consensus protocol, the security loop gains strong structural resilience. At blockchain level, this mechanism tolerates the active compromise of up to $1/3$ of validators while preserving ledger safety and finality under partial network compromise. This guarantee is distinct from IDS-side voting, which is the mechanism used to validate the semantic correctness of observations before they affect the consolidated monitored state. The two layers compose into a two-threshold fault model: an attacker must subvert either a majority of the heterogeneous IDSs observing a variable or a supermajority of the consensus validators, a requirement that makes the compromise of the autonomic manager highly improbable.

Generally, the system is designed to reduce the attack surface as much as possible by treating actuators as black boxes (blind). The actuators receive signed mitigation commands directly from the blockchain state, rather than querying the state of potentially compromised application containers directly. This enforces a strict unidirectional data flow: IDSs push observations to the consensus plane, the blockchain validates and commits the state transition, and the resulting response is pushed to the actuators: so that an actuator only ever executes commands that have passed through BFT consensus.

A final, often overlooked, property is \emph{deterministic finality}. A software control loop cannot act on decisions that might subsequently be rolled back, so the underlying consensus must commit blocks that are final and immutable rather than probabilistically finalized. This rules out the public, permissionless blockchain model and its Proof-of-Work/Proof-of-Stake machinery, with their tokenomics, high latency, and resource consumption, and it motivates the choice of a permissioned environment where participants are pre-vetted, identities are static, and consensus can be simplified to a BFT protocol with immediate finality.

% : : : : : : : : : : : : : : : : : : : : : : : : : : 
\section{The Technology Choice Is Not Neutral}
\label{sec:intro_tech_choice}
% : : : : : : : : : : : : : : : : : : : : : : : : : : 

While the theoretical integration of blockchain into the SPS architecture provides outstanding security guarantees, the practical choice of the underlying blockchain technology is not neutral. A blockchain-based autonomic manager has unique requirements: it must handle sparse, highly event-driven alert traffic with exceptionally low latency and high determinism, and it must absorb bursty, redundant alert storms without collapsing into unbounded latency or silently dropping observations. Conversely, many features native to public blockchains: virtual-machine-level gas metering, token economics, transaction fees, dynamic peer discovery, and an archival storage model: add unnecessary complexity, expand the trusted computing base (TCB), and increase the attack surface without contributing to the control loop.

The natural baseline choice for permissioned enterprise setups is GoQuorum, a fork of Ethereum that retains the Ethereum Virtual Machine (EVM) and provides BFT consensus through IBFT/QBFT. However, this thesis argues that general-purpose smart-contract platforms incur severe, non-intrinsic performance overheads due to interpreted virtual machine bytecode execution, rigid transaction fee mechanisms, and fixed-interval block production: every alert pays RLP serialization, EVM bytecode interpretation, and instruction-granularity gas metering that execute unconditionally even though fees are zero, and the fixed block period caps how quickly the mempool can drain. As an alternative, this work explores CometBFT (the successor of Tendermint), a lower-level BFT consensus engine that keeps consensus separate from the application and supports highly optimized, application-specific distributed state machines through the Application BlockChain Interface (ABCI). In CometBFT the SPS state machine can be implemented as a native function call, empty blocks can be disabled so that the first transaction is proposed immediately, and the runtime cost of a state transition becomes proportional to its intrinsic logic rather than to virtual-machine scaffolding.

The comparison is framed through a deliberately requirement-driven lens: a feature that is not needed by the SPS manager is best \emph{absent}, because any additional subsystem increases configuration complexity, wastes computational power, and creates new failure modes or attack vectors that must themselves be secured. Under this lens, the absence of non-required functionality in CometBFT is itself a form of alignment with the SPS use case.

% : : : : : : : : : : : : : : : : : : : : : : : : : : 
\section{Objectives and Contributions of the Thesis}
\label{sec:intro_objectives}
% : : : : : : : : : : : : : : : : : : : : : : : : : : 

The primary objective of this thesis is to present a comprehensive qualitative and empirical comparison between a conventional general-purpose platform (GoQuorum) and a requirement-driven, application-specific alternative (CometBFT) to determine which architecture best supports the performance and cyber-resilience demands of modern Self-Protecting Systems. The comparison is requirement-driven rather than feature-driven: it does not reward a platform for offering additional services unless those services directly improve the SPS manager's correctness, latency, or operational resilience.

Specifically, the core contributions of this thesis are as follows:
\begin{enumerate}
    \item \textbf{Requirement Formulation:} We define and classify a comprehensive set of functional and non-functional requirements (required, desirable, and not needed) for blockchain technology in the specific context of an autonomic SPS loop, derived from the architecture and aligned with the fixed trust boundary of a permissioned deployment.
    \item \textbf{Architectural Analysis:} We conduct a thorough qualitative comparison of GoQuorum and CometBFT, highlighting the architectural mismatch of general-purpose virtual machines versus application-specific state machines, and reading each platform's design choices onto the SPS requirement set.
    \item \textbf{Prototype Implementation:} We design and implement a realistic, fully containerized SPS prototype emulated via Kathará on top of Docker. The prototype integrates three heterogeneous IDS engines (Snort, Suricata, and Zeek) protecting a vulnerable web application (OWASP Juice Shop), and runs interchangeably on both blockchain backends so that any measured difference is attributable to the consensus substrate.
    \item \textbf{Deduplication Middleware:} We propose and develop a custom deduplication middleware layer for the member nodes. This layer filters repetitive alerts before they reach the mempool, collapsing temporally adjacent observations that request the same state transition into a single on-chain transaction and keeping the transaction count upper-bounded per block.
    \item \textbf{Empirical Evaluation:} We perform extensive automated benchmarks under controlled attack conditions, comparing GoQuorum and CometBFT across key performance indicators: end-to-end response latency, transaction throughput under heavy load, and the compression ratio achieved by the deduplication middleware: and relate the measured behaviour back to the qualitative predictions of the requirement analysis.
\end{enumerate}

% : : : : : : : : : : : : : : : : : : : : : : : : : : 
\section{Structure of the Thesis}
\label{sec:intro_structure}
% : : : : : : : : : : : : : : : : : : : : : : : : : : 

To thoroughly explore the theoretical background, architectural design, and empirical evaluation of the proposed systems, the remainder of this thesis is structured as follows:

\begin{description}
    \item[Chapter 2 -- Technologies Setting and Related Work:] Introduces the foundational concepts of autonomic computing and the MAPE-K cycle, the structure of permissioned blockchain environments, the family of distributed consensus algorithms from Crash Fault Tolerance to Byzantine Fault Tolerance, and reviews the state of the art in autonomic manager security, from traditional hardening to Trusted Execution Environments and distributed ledger technologies.
    \item[Chapter 3 -- A Cyber-Resilient Architectural Model:] Details the technology-agnostic architecture of the Self-Protecting System, formalizing the system decomposition into detection, consensus-and-decision, and control planes, the replicated state machine that governs the autonomic loop, the end-to-end protection workflow, and the safety, liveness, and threat-model properties that the architecture guarantees.
    \item[Chapter 4 -- Requirement Alignment: Quorum vs. CometBFT:] Formally identifies the blockchain properties strictly required, desirable, and not needed for SPS applications, and theoretically compares Quorum's general-purpose execution environment with CometBFT's application-specific design through a requirement-alignment table.
    \item[Chapter 5 -- Implementation Details:] Describes the prototype implementation, the containerized network topology emulated via Kathará, the realization of the replicated decision state machine as both a Solidity contract and a Rust ABCI application, the end-to-end alert response loop, and the custom deduplication middleware developed for the evaluation.
    \item[Chapter 6 -- Experiments and Results:] Analyzes the performance of the two blockchain backends under three research questions: end-to-end latency under low traffic, sustained throughput under increasing load, and the effectiveness of the deduplication middleware: presenting metrics on response time, transaction throughput, and the compression of redundant attack traffic.
    \item[Chapter 7 -- Conclusions and Future Work:] Summarizes the findings of the comparative analysis, acknowledges the limitations of the study in terms of experimental setting, threat-model scope, and engineering maturity, and outlines directions for future research including physical cluster deployments, richer attack scenarios, sharding, and formal verification of the state transition function.
\end{description}
